\section{Basic Concepts of Modules}

Let \(R\) be a ring and \(\para{M, +, 0}\) be an abelian group.

\begin{definition}[\(R\)-Module]
    \label{def:r-module}%
    A quadruple \(\para{M, \cdot}\) is an \emph{\(R\)-module} if the operation
    \[
        \function{\cdot}{R \times M}{M}{\para{r, m}}{rm}
    \]
    satisfies the following.

    \begin{enumerate}
        \item \textit{Distributivity over Scalar.} \(\para{r_1 + r_2} m = r_1 m + r_2 m\).
        \item \textit{Distributivity over `Vector'.} \(r \para{m_1 + m_2} = r m_1 + r m_2\).
        \item \textit{`Associativity' in Scalar Multiplication.} \(r_1 \cdot \para{r_2 m} = \para{r_1 \cdot r_2} m\).
        \item \textit{Unitary Condition.} \(1_{R} \cdot m = m\).
    \end{enumerate}
\end{definition}

\begin{examples}[Examples of \(R\)-Modules]
    \label{ex:examples-r-modules}%
    \begin{itemize}
        \item If \(R = k\) is a field, then an \(R\)-module \(M\) is the same thing as a \(k\)-vector space.
        \item Any abelian group \(G\) is a \(\ZZ\)-module, by the operation
        \[
            n \cdot g = g^n.
        \]
        \item Take any ring \(R\). The product \(R^m = R \times \cdots \times R\) is an \(R\)-module, via the action
        \[
            \function{\cdot}{R \times R^m}{R^m}{\para{r, \para{r_1, \ldots, r_m}}}{\para{r \cdot r_1, \ldots, r \cdot r_m}.}
        \]

        [Note \(R^m\) is also equipped with a ring structure, but we do not use it here.]

        \item If \(I \idealsub R\) is an ideal, then \(I\) is an \(R\)-module, via the multiplication action
        \[
            \function{\cdot}{R \times I}{I}{\para{r, x}}{rx.}
        \]

        Similarly, \(R / I\) is also an \(R\)-module, via the multiplication action
        \[
            \function{\cdot}{R \times R / I}{R / I}{\para{r, x + I}}{\para{r + I} \cdot \para{x + I}.}
        \]
    \end{itemize}
\end{examples}

\begin{example*}[\(k\)-Vector Space as a \(k\brac{x}\)-Module]
    Let \(V\) be a \(k\)-vector space, and we fix \(\alpha \in \End\para{V}\) be an endomorphism of \(V\), i.e., a linear map \(\functiondomain{\alpha}{V}{V}\). We can make \(V\) into a \(k\brac{x}\)-module as follows
    \[
        \function{\cdot_{\alpha}}{k\brac{x} \times V}{V}{\para{f, \vm{v}}}{\para{f\para{\alpha} \vm{v}}.}
    \]
\end{example*}

\begin{example*}[Ring Homomorphism and Modules]
    Take a ring homomorphism \(\functiondomain{\phi}{R}{S}\). This makes \(S\) into an \(R\)-module, with
    \[
        \function{\cdot}{R \times S}{S}{\para{r, s}}{\phi\para{r} \cdot s.}
    \]

    For example, the quotient map \(\functiondomain{q}{\CC\brac{x}}{\CC\brac{x} / x^2}\) makes \(\CC\brac{x} / x^2\) a \(\CC\brac{x}\)-module, and the inclusion map \(\functiondomain{i}{\CC\brac{x}}{\CC\brac{x, y}}\) makes \(\CC\brac{x, y}\) a \(\CC\brac{x}\) module.
\end{example*}

\begin{definition}[Submodule]
    \label{def:submodule}%
    Let \(M\) be an \(R\)-module. A subgroup \(N\) of \(M\) is called a \emph{submodule}, if for any \(n \in N\) and \(r \in R\), then \(r \cdot n \in N\). We denote this as \(N \submodule M\).

    Given an element \(m \in M\), the \emph{submodule generated by \(m\)}, is
    \[
        Rm = \set{rm}{r \in R} \submodule M.
    \]
\end{definition}

\begin{example*}[Non-Example of Submodule]
    Take \(R = M = \QQ\), and \(N = \ZZ \subseteq \QQ\) is not a \(\QQ\)-submodule, since \(\ZZ\) is not a \(\QQ\)-vector subspace of \(\QQ\).
\end{example*}

\begin{remark}
    If \(R\) is a field, an \(R\)-submodule is the same as an \(R\)-vector subspace.
\end{remark}

\begin{definition}[Quotient by Submodule]
    \label{def:quotient-submodule}%
    Let \(N \submodule M\) be an \(R\)-submodule. The \emph{quotient} \(R\)-module is the abelian group \(M / N\), with the \(R\)-module structure
    \[
        \function{\cdot}{R \times M / N}{M / N}{\para{r, m + N}}{\para{rm} + N.}
    \] 
\end{definition}

\begin{definition}[Module Homomorphism and Isomorphism]
    \label{def:homomorphism-isomorphism-module}%
    Let \(M, N\) be \(R\)-modules. An \emph{\(R\)-module homomorphism} is a group homomorphism \(\functiondomain{\phi}{M}{N}\), such that
    \[
        \phi \para{r \cdot m} = r \cdot \phi\para{m}
    \]
    for all \(r \in R, m \in M\).

    A bijective \(R\)-module homomorphism is called an \emph{\(R\)-module isomorphism}.
\end{definition}

\begin{theorem}[First Isomorphism Theorem for Modules]
    \label{thm:first-isomorphism-theorem-modules}%
    Let \(\functiondomain{\phi}{M}{N}\) be a homomorphism of \(R\)-modules. Then
    \[
        \ker\para{\phi} = \set{m \in M}{\phi\para{m} = 0}
    \]
    is an \(R\)-submodule of \(M\), and the image
    \[
        \img\para{\phi} = \set{\phi\para{m} \in N}{m \in M}
    \]
    is an \(R\)-submodule of \(N\), and we have an isomorphism
    \[
        \function{\psi}{M / \ker \phi}{\img \phi}{m + \ker \phi}{\phi\para{m}.}
    \]
\end{theorem}

\begin{theorem}[Second Isomorphism Theorem for Modules]
    \label{thm:second-isomorphism-theorem-modules}%
    Let \(K, L \submodule M\). Then the set
    \[
        K + L = \set{k + l}{k \in K, l \in L}
    \]
    is a submodule of \(M\). Furthermore, there is an isomorphism
    \[
        \para{K + L} / K \isomorphic L / \para{K \cap L}.
    \]
\end{theorem}

\begin{theorem}[Third Isomorphism Theorem for Modules]
    \label{thm:third-isomorphism-theorem-modules}%
    Let \(N \submodule L \submodule M\) be submodules of \(M\). Then
    \[
        M / L \isomorphic \para{M / N} / \para{L / N}.
    \]
\end{theorem}

There is also a correspondence result. Let \(N \submodule M\), then the sets
\[
    \brce{\text{submodules of }M / N} \leftrightarrow \brce{\text{submodules of }M\text{ that contain }N}
\]
are bijective.

\begin{definition}[Annihilator]
    \label{def:annihiliator}%
    Let \(M\) be an \(R\)-submodule. For \(m \in M\), its \emph{annihilator} is
    \[
        \Ann\para{m} = \set{r \in R}{r \cdot m = 0}.
    \]

    More generally, for \(S \subseteq M\), the \emph{annihilator} of \(S\) is
    \[
        \Ann\para{S} = \set{r \in R}{\forall s \in S: r \cdot s = 0}.
    \]
\end{definition}

The annihilator of the whole module is
\[
    \Ann\para{M} = \bigcap_{m \in M} \Ann\para{m}
\]
is similar to the idea of the kernel of a group action -- since we may think of modules as an `action of a ring on an abelian group'.

When \(R\) is a field, \(M\) is an \(R\)-vector space, so we have
\[
    \Ann\para{m} = \brce{0}
\]
for any \(m \in M\). This notion is useless in the theory of vector spaces, but not in general for modules.

\begin{example*}[Annihilator]
    Take \(R = \ZZ\), and \(M = \ZZ / n\ZZ\) with the usual scaling. The annihilator of this is interesting, for example, for \(1 \in M = \ZZ / n \ZZ\), we have \(\Ann\para{1} = n\ZZ\).
\end{example*}

\begin{proposition}[Annihilator is an Ideal]
    \label{prop:annihilator-ideal}%
    The annihilator \(\Ann\para{S}\) is an ideal in \(R\).
\end{proposition}

\begin{proposition}[Isomorphism between Annihilator and Submodule]
    \label{prop:isomorphism-annihilator-submodule}%
    For a fixed \(m \in M\), there is an isomorphism
    \[
        R / \Ann\para{m} \isomorphic Rm.
    \]
\end{proposition}

\begin{proof}
    Consider the function
    \[
        \function{\phi_m}{R}{M}{r}{rm}
    \]
    for some fixed \(m \in M\). This is an \(R\)-module homomorphism (note \(R\) is trivially an \(R\)-module). Applying \autoref{thm:first-isomorphism-theorem-modules} (First Isomorphism Theorem of Modules) gives the result.
\end{proof}

\begin{remark}
    If \(M\) were a vector space over a field \(R\), then the submodule (subspace) generated by \(m \in M\) is
    \[
        Rm = \spn\brce{m}
    \]
    is a (one-dimensional) line -- but this intuition is not true for rings, because there are non-trivial annihilators!
\end{remark}

\begin{definition}[Finitely Generated Module]
    \label{def:finitely-generated-module}%
    An \(R\)-module \(M\) is \emph{finitely generated} if there exists some \(m_1, \ldots, m_k \in M\), such that
    \[
        M = R m_1 + \cdots + R m_k = \set{r_1 m_1 + \cdots + r_k m_k}{r_i \in R}.
    \]
\end{definition}

This is a similar idea to a finite dimensional vector space -- but we are not saying that this is a basis -- it is just a spanning set!

\begin{lemma}[Condition on Finitely Generated Modules]
    \label{lem:finitely-generated-modules-surjective-homomorphism}%
    An \(R\)-module \(M\) is finitely generated, if and only if there exists a surjective \(R\)-module homomorphism \(R^k \to M\) for some \(k\).
\end{lemma}

In other words, \(M\) is finitely generated if it is a quotient of \(R^k\) for some \(k\).

\begin{proof}
    If \(M\) is finitely generated, then we can write
    \[
        M = R m_1 + \cdots + R m_k
    \]
    as above. Then we define the homomorphism
    \[
        \function{\phi}{R^k}{M}{\para{r_1, \ldots, r_k}}{r_1 m_1 + \cdots + r_k m_k}
    \]
    which is clearly surjective.

    Conversely, given some surjective homomorphism \(\functiondomain{\phi}{R^k}{M}\), let \(e_i = \para{0, \ldots, 0, 1, 0, \ldots, 0} \in R^k\) with \(1\) in the \(i\)th component.

    Define \(m_i = \phi\para{e_i}\). It is straightforward to check that \(M = R m_1 + \cdots + R m_k\).
\end{proof}

The modules \(R^k\) are often called \emph{free \(R\)-modules}.

\begin{corollary}[Quotient of Finitely Generated Module is Finitely Generated]
    \label{cor:quotient-finitely-generated}%
    Let \(M\) be finitely generated. If \(N \submodule M\) is a submodule, then the quotient \(M / N\) is finitely generated as well.
\end{corollary}

\begin{proof}
    Take surjective \(R^k \to M\). The composition with the quotient map \(M \to M / N\) is a composition of surjective maps, hence surjective.
\end{proof}

\begin{remark}
    A submodule of a finitely generated module need not be finitely generated. Take
    \[
        R = \CC\brac{x_1, x_2, \ldots}
    \]
    be a polynomial ring over countably many variables, and let \(M = R\). This is clearly finitely generated.

    However, \(I \submodule M\) to be the ideal generated by
    \[
        I = \para{x_1, x_2, \ldots},
    \]
    which consists of all the polynomials without constant terms. This is not finitely generated.
\end{remark}

